\scanpage{108}
de \deleted{niveau} poids $\le 2q$, de niveau $\le q$. Donc il en est ainsi de $E^{rs}_2$ si $p+q=i$, $p\ne 0$, \inserted{et ce terme est de \deleted{\unsure{poids, de}} poids $\le 2(i-1)<2i$, de niveau $\le i-1<i$} Reste à examiner le terme $E^{0,i}_2\subset H^i(X')$ ; comme $X'$ est lisse, il est donc effectif, de poids $\le 2i$, de niveau $\le i$. On gagne\unsure{!}, si on note que si $i>n$, on a $H^i(X)=0$ ($X$ affine !).

\deleted{Si $X$ \unsure{propre ;|provenait ;|fortiori} $H^i(X)$ est de niveau $\le n=\dim X$. Il suffit \unsure{de l'éprouver pour les} $E_2^{pq}$. \unsure{Il suffit d'utiliser un recouvt}}

\textcircled{B} \uline{Coefficients constants tordus sur $X$ lisse} \inserted{séparé}, $F$ \inserted{effectif} de poids et \inserted{niveau $j$.}

\textcircled{1} On suppose que les coefficients envisagés proviennent d'un motif lisse sur $X$, et \inserted{\unsure{par dévissage (possible loc ?)|par dévissages (possible loc ?)}} on va regarder le cas où ce dernier est un facteur direct d'un $R^jf_*(\mathbb Q_{\ell Y})$, avec $f:\deleted{\unsure{X}}Y\to X$ morphisme projectif et lisse. \inserted{On peut si on veut supposer $f$ de dim relative $j$.} On est donc ramené au cas où $F=R^jf_*(Q_{\ell Y})$. \inserted{On} Considérons la suite spectrale
\[
 H^*(Y)\deleted{\unsure{\Leftarrow}}\Longleftarrow \inserted{E^{pq}_2=}H^p(X,R^q f_*(Y))
\]

\uuline{D'après Deligne,} cette suite spectrale dégénère. Donc $H^i(X,R^jf_*(Q_{\ell Y}))$ \inserted{et par suite $H^i(X,F)$} apparait comme un \uline{morceau} de $H^{i+j}(Y)$. On trouve que l'on a un \deleted{module} \uline{effectif}, de poids \deleted{\unsure{\#}} compris entre \deleted{\unsure{$0$}} \deleted{$2(n+j)$ et $2(i+j)$} \uline{$i+j$} et $2(n+j)=\underline{2n+2j}$ (où $n=\dim X$) ; les niveaux sont resp. $\le(2n-i+j),\ldots,0$.
\margin{peut-on remplacer $2n+2j$ par $2n+j$ ?}

Si $i\le n$, il n'y a de poids \unsure{qui ont|qu'en} les dimensions $i+j,i+j+1,\ldots,\bigl(2(i+j)\bigr)$, de \deleted{poids} \inserted{niveaux} respectifs \deleted{\unsure{$2i,\ldots,i+j$|$\le 0,\ldots,i+j$}.} \inserted{$i+j,i+j-1,\ldots,0$}
\margin{peut-on remplacer $2i+2j$ par $2i+j$ ? \deleted{\unsure{Il doit faire ça pour|Idiot \illegible{3}}} !}

Résultats duels pour $H^i_!(F)$, il est effectif à poids $\le i+j$, et \deleted{à niveaux} si $i\ge n$, il est à niveaux croiss. de de $0$ à \deleted{$2n-i+j$} $2(n+j)-(i+j)=2n-i+j$. \inserted{pour des poids croissants de}
\margin{les poids sont-ils $\ge j$ \deleted{\unsure{resp.}}}

Si $X$ est \uline{propre}, on trouve que $H^i(X,F)$ est pur de poids $i+j$ (et semi-simple ici\ldots, !
% archive sheet number: 107
